% GENERATED FILE. Edit canonical lessons, then run npm run generate:books.
% canonical-source-hash: fnv1a64:6a6469a9abe29ccd
% canonical-lessons: BN-C323-nayok, BN-C323-kholnayok, BN-C323-bhumika, BN-C323-sonnyasi, BN-C323-chhodmobesh

\chapter{Heroes and Villains}
\label{ch:bn-heroes-and-villains}
\hlchaptermodality{323}

\begin{chapteropening}
\textbf{By the end of this chapter:} \emph{I can talk about the people in a story and the parts they play.}

Complete the last lesson of chapter 323: I can talk about the people in a story and the parts they play.
\end{chapteropening}

\section[\texorpdfstring{\bn{নায়ক} (nāyôk)}{নায়ক (nāyôk)}]{\bn{নায়ক} (nāyôk) — the hero of a story; a protagonist}

\label{lesson:BN-C323-nayok}



\begin{quote}
Before the new one: say the Bengali for a curse, then the Bengali for a ghost.
\end{quote}

\subsection*{You'll want to know: \bn{নায়ক}}
\textbf{\bn{নায়ক}} — \emph{nāyôk} — \textquotedblleft{}the hero of a story; a protagonist\textquotedblright{}.

\begin{quote}
\bn{গল্পের} \bn{নায়ক} \bn{এক} \bn{গরিব} \bn{জেলে।}
\end{quote}

\begin{quote}
\emph{gôlper nāyôk ek gôrib jele.}
\end{quote}

\begin{quote}
The hero of the story is a poor fisherman.
\end{quote}

\begin{cousinweb}
\textbf{\bn{নায়ক}} (\emph{nāyôk}) is Sanskrit \emph{nāyaka}, 'a guide, a leader, a chief' (Monier-Williams); in plays and stories it is the leading man.
\end{cousinweb}

\begin{grammarlens}[title={What you've built}]
One more word for telling what happened.
\end{grammarlens}

\subsection*{Your turn}
\begin{itemize}
\raggedright
  \item \emph{Say it:} \emph{nāyôk}
  \item \emph{Say it:} \emph{nāyôk}, once more
  \item \emph{Say it:} \emph{gôlper nāyôk ek gôrib jele.} — The hero of the story is a poor fisherman.
  \item \emph{Recall:} say the Bengali for a curse, then the Bengali for a ghost, then say \emph{nāyôk} again
\end{itemize}

\begin{tcolorbox}[breakable,colback=teal!4,colframe=teal!35!black,title={Before you move on}]
What does \bn{নায়ক} mean? (\textquotedblleft{}The hero of a story; a protagonist\textquotedblright{}.) Say it once more.
\end{tcolorbox}
\section[\texorpdfstring{\bn{খলনায়ক} (khôlnāyôk)}{খলনায়ক (khôlnāyôk)}]{\bn{খলনায়ক} (khôlnāyôk) — a villain}

\label{lesson:BN-C323-kholnayok}



\begin{quote}
Before the new one: say the Bengali for a ghost, then the Bengali for the hero of a story.
\end{quote}

\subsection*{You'll want to know: \bn{খলনায়ক}}
\textbf{\bn{খলনায়ক}} — \emph{khôlnāyôk} — \textquotedblleft{}a villain\textquotedblright{}.

\begin{quote}
\bn{শেষে} \bn{পুলিশ} \bn{খলনায়ককে} \bn{গ্রেফতার} \bn{করল।}
\end{quote}

\begin{quote}
\emph{sheshe pulish khôlnāyôkke grephtār kôrlô.}
\end{quote}

\begin{quote}
In the end the police arrested the villain.
\end{quote}

\begin{cousinweb}
\textbf{\bn{খলনায়ক}} (\emph{khôlnāyôk}) joins Sanskrit \emph{khala}, 'a mischievous man, a rogue' (Monier-Williams), to \textbf{\bn{নায়ক}}: the rogue who leads the other side of the story.
\end{cousinweb}

\begin{grammarlens}[title={What you've built}]
One more word for telling what happened.
\end{grammarlens}

\subsection*{Your turn}
\begin{itemize}
\raggedright
  \item \emph{Say it:} \emph{khôlnāyôk}
  \item \emph{Say it:} \emph{khôlnāyôk}, once more
  \item \emph{Say it:} \emph{sheshe pulish khôlnāyôkke grephtār kôrlô.} — In the end the police arrested the villain.
  \item \emph{Recall:} say the Bengali for a ghost, then the Bengali for the hero of a story, then say \emph{khôlnāyôk} again
\end{itemize}

\begin{tcolorbox}[breakable,colback=teal!4,colframe=teal!35!black,title={Before you move on}]
What does \bn{খলনায়ক} mean? (\textquotedblleft{}A villain\textquotedblright{}.) Say it once more.
\end{tcolorbox}
\section[\texorpdfstring{\bn{ভূমিকা} (bhūmikā)}{ভূমিকা (bhūmikā)}]{\bn{ভূমিকা} (bhūmikā) — a role; a preface}

\label{lesson:BN-C323-bhumika}



\begin{quote}
Before the new one: say the Bengali for the hero of a story, then the Bengali for a villain.
\end{quote}

\subsection*{You'll want to know: \bn{ভূমিকা}}
\textbf{\bn{ভূমিকা}} — \emph{bhūmikā} — \textquotedblleft{}a role; a preface\textquotedblright{}.

\begin{quote}
\bn{নাটকে} \bn{মীরার} \bn{ভূমিকা} \bn{ছিল} \bn{এক} \bn{ডাইনির।}
\end{quote}

\begin{quote}
\emph{nāṭôke mīrār bhūmikā chhilô ek ḍāinir.}
\end{quote}

\begin{quote}
In the play Mira's part was a witch.
\end{quote}

\begin{cousinweb}
\textbf{\bn{ভূমিকা}} (\emph{bhūmikā}) is Sanskrit \emph{bhūmikā}, 'ground, a place', which Monier-Williams also gives as an actor's part and as the preface of a book; Bengali keeps both.
\end{cousinweb}

\begin{grammarlens}[title={What you've built}]
One more word for telling what happened.
\end{grammarlens}

\subsection*{Your turn}
\begin{itemize}
\raggedright
  \item \emph{Say it:} \emph{bhūmikā}
  \item \emph{Say it:} \emph{bhūmikā}, once more
  \item \emph{Say it:} \emph{nāṭôke mīrār bhūmikā chhilô ek ḍāinir.} — In the play Mira's part was a witch.
  \item \emph{Recall:} say the Bengali for the hero of a story, then the Bengali for a villain, then say \emph{bhūmikā} again
\end{itemize}

\begin{tcolorbox}[breakable,colback=teal!4,colframe=teal!35!black,title={Before you move on}]
What does \bn{ভূমিকা} mean? (\textquotedblleft{}A role; a preface\textquotedblright{}.) Say it once more.
\end{tcolorbox}
\section[\texorpdfstring{\bn{সন্ন্যাসী} (sônnyāsī)}{সন্ন্যাসী (sônnyāsī)}]{\bn{সন্ন্যাসী} (sônnyāsī) — a holy man; an ascetic}

\label{lesson:BN-C323-sonnyasi}



\begin{quote}
Before the new one: say the Bengali for a villain, then the Bengali for a role.
\end{quote}

\subsection*{You'll want to know: \bn{সন্ন্যাসী}}
\textbf{\bn{সন্ন্যাসী}} — \emph{sônnyāsī} — \textquotedblleft{}a holy man; an ascetic\textquotedblright{}.

\begin{quote}
\bn{এক} \bn{সন্ন্যাসী} \bn{রাজকুমারকে} \bn{অভিশাপ} \bn{থেকে} \bn{বাঁচালেন।}
\end{quote}

\begin{quote}
\emph{ek sônnyāsī rājkumārke ôbhishāp theke bā̃chālen.}
\end{quote}

\begin{quote}
A holy man saved the prince from the curse.
\end{quote}

\begin{cousinweb}
\textbf{\bn{সন্ন্যাসী}} (\emph{sônnyāsī}) is Sanskrit \emph{saṃnyāsin}, 'laying aside, renouncing' (Monier-Williams): one who has given up the world.
\end{cousinweb}

\begin{grammarlens}[title={What you've built}]
One more word for telling what happened.
\end{grammarlens}

\subsection*{Your turn}
\begin{itemize}
\raggedright
  \item \emph{Say it:} \emph{sônnyāsī}
  \item \emph{Say it:} \emph{sônnyāsī}, once more
  \item \emph{Say it:} \emph{ek sônnyāsī rājkumārke ôbhishāp theke bā̃chālen.} — A holy man saved the prince from the curse.
  \item \emph{Recall:} say the Bengali for a villain, then the Bengali for a role, then say \emph{sônnyāsī} again
\end{itemize}

\begin{tcolorbox}[breakable,colback=teal!4,colframe=teal!35!black,title={Before you move on}]
What does \bn{সন্ন্যাসী} mean? (\textquotedblleft{}A holy man; an ascetic\textquotedblright{}.) Say it once more.
\end{tcolorbox}
\section[\texorpdfstring{\bn{ছদ্মবেশ} (chhôdmôbesh)}{ছদ্মবেশ (chhôdmôbesh)}]{\bn{ছদ্মবেশ} (chhôdmôbesh) — a disguise}

\label{lesson:BN-C323-chhodmobesh}



\begin{quote}
Before the new one: say the Bengali for a role, then the Bengali for a holy man.
\end{quote}

\subsection*{You'll want to know: \bn{ছদ্মবেশ}}
\textbf{\bn{ছদ্মবেশ}} — \emph{chhôdmôbesh} — \textquotedblleft{}a disguise\textquotedblright{}.

\begin{quote}
\bn{খলনায়ক} \bn{সন্ন্যাসীর} \bn{ছদ্মবেশে} \bn{প্রাসাদে} \bn{ঢুকল।}
\end{quote}

\begin{quote}
\emph{khôlnāyôk sônnyāsīr chhôdmôbeshe prāsāde ḍhuklô.}
\end{quote}

\begin{quote}
The villain got into the palace disguised as a holy man.
\end{quote}

\begin{cousinweb}
\textbf{\bn{ছদ্মবেশ}} (\emph{chhôdmôbesh}) joins Sanskrit \emph{chadman}, 'a covering, a deceptive dress' (Monier-Williams), to \emph{veśa}, 'dress': false clothes.
\end{cousinweb}

\begin{grammarlens}[title={What you've built}]
One more word for telling what happened.
\end{grammarlens}

\subsection*{Your turn}
\begin{itemize}
\raggedright
  \item \emph{Say it:} \emph{chhôdmôbesh}
  \item \emph{Say it:} \emph{chhôdmôbesh}, once more
  \item \emph{Say it:} \emph{khôlnāyôk sônnyāsīr chhôdmôbeshe prāsāde ḍhuklô.} — The villain got into the palace disguised as a holy man.
  \item \emph{Recall:} say the Bengali for a role, then the Bengali for a holy man, then say \emph{chhôdmôbesh} again
\end{itemize}

\begin{tcolorbox}[breakable,colback=teal!4,colframe=teal!35!black,title={Before you move on}]
What does \bn{ছদ্মবেশ} mean? (\textquotedblleft{}A disguise\textquotedblright{}.) Say it once more.
\end{tcolorbox}
